\section{FORMAL RESTATEMENT}
\textbf{Erd\H{o}s \#378; reconstruction with endpoint correction, 2026-09-24.}
Let
\[
S(n)=\#\{1\leq j\leq n-1:\ \binom nj\text{ is squarefree}\}.
\]
For each fixed positive integer $r$, establish existence of a positive natural density for $\{n:S(n)\geq r\}$.

\section{QUICK LITERATURE AND CONTEXT CHECK}
Granville and Ramar\'e proved this in \emph{Explicit bounds on exponential sums and the scarcity of squarefree binomial coefficients}. Their Theorem 5 counts the whole Pascal row, including $\binom n0=\binom nn=1$. Thus their $2m+2$ squarefree entries correspond to $2m$ \emph{internal} entries. Applying $2m+2$ directly to the problem's internal count would give an incorrect index shift.

\section{OBJECTIVE AND ATTACK PLAN}
Use the published distribution theorem with its exact endpoint convention. Prove the parity reduction, derive a finite complementary sum for each requested density, and give an elementary positive lower bound for the first two thresholds.

\section{WORK}
\textbf{External inputs, stated precisely.} Theorem 5 of Granville--Ramar\'e gives a natural density $\eta_m$ for the integers $n$ whose whole Pascal row has exactly $2m+2$ squarefree entries, for every integer $m\geq0$. Moreover $\eta_m>0$ for every $m\geq1$. Their Theorem 1 states that $\binom{2a}{a}$ is not squarefree for $a>4$. These two results are used as published inputs; the exponential-sum and sieve estimates establishing them are not reproduced here.

\textbf{Why the number of internal entries is even.} The involution $j\mapsto n-j$ preserves $\binom nj$ and pairs all internal indices except possibly $j=n/2$. For odd $n$ there is no fixed index. For even $n>8$ the fixed central coefficient is not squarefree by the second input. Hence $S(n)$ is even for every $n>8$.

Since the two endpoints always contribute two squarefree entries, the set
\[
E_m=\{n:S(n)=2m\}
\]
has density $\eta_m$. For $n>8$, failure of $S(n)\geq r$ is exactly membership in one of the finitely many disjoint sets $E_m$ with $2m<r$. Therefore the desired density exists and is
\[
\boxed{\ \delta_r=1-\sum_{m=0}^{\lceil r/2\rceil-1}\eta_m\ }.       \tag{1}
\]
Only a finite union is involved in this deduction, so no interchange of an infinite sum and a density limit is being assumed. If $m\geq\max(1,\lceil r/2\rceil)$, then $E_m\subseteq\{n:S(n)\geq r\}$ and $\eta_m>0$; hence $\delta_r>0$. Also (1) gives $\delta_{2h-1}=\delta_{2h}$ for $h\geq1$, reflecting the parity pairing.

\textbf{An elementary check on the first thresholds.} For squarefree $n\geq3$, the two distinct entries at $j=1,n-1$ both equal $n$, so $S(n)\geq2$. The identity
\[
\mathbf1_{\text{squarefree}}(n)=\sum_{d^2\mid n}\mu(d)
\]
gives
\[
\#\{n\leq x:n\text{ squarefree}\}
=\sum_{d\leq\sqrt x}\mu(d)\lfloor x/d^2\rfloor
=\frac{x}{\zeta(2)}+O(\sqrt x).
\]
Here the rounding errors total $O(\sqrt x)$ and the omitted absolutely convergent series tail contributes $O(\sqrt x)$. Thus
\[
\delta_1=\delta_2\geq\frac6{\pi^2}.
\]
This independently verifies positivity at the first two thresholds, but does not identify their exact density.

\section{RESULT AND LIMITATIONS}
\textbf{Attempt status: solved relative to the two explicitly stated Granville--Ramar\'e theorems.} Formula (1), positivity for every fixed $r$, and the endpoint and parity conventions are fully justified. The deep published distribution theorem remains an external dependency; no new distribution theorem or numerical values of $\eta_m$ are claimed.

\section{SOURCES}
\url{https://www.erdosproblems.com/378}; Granville and Ramar\'e, Theorems 1 and 5 (printed pages 1 and 3): \url{https://dms.umontreal.ca/~andrew/PDF/ramare.pdf}.

\textbf{Completion rate estimate: 100\% relative to the stated external inputs.}
